\chapter{Introduction}

\section{Purpose}

The Narrative Bible defines the canonical lore, history, cultures, factions, locations, and storytelling principles of \textit{ECHO: The Living World}. It serves as the authoritative reference for writers, designers, artists, voice actors, and developers throughout the project's lifecycle.

While the Game Design Document explains what the player experiences and the Technical Design Document explains how systems are implemented, the Narrative Bible establishes the fictional universe in which those experiences take place.

\section{Objectives}

The primary objectives of this document are:

\begin{itemize}
    \item Establish the official lore of Aethera.
    \item Maintain consistency across all narrative content.
    \item Define historical events and world-building.
    \item Document factions, cultures, and civilizations.
    \item Support future expansions and sequels.
\end{itemize}

\section{Scope}

This document covers the world's history, geography, civilizations, political powers, mythology, major characters, narrative rules, and future storytelling direction.

Gameplay mechanics and technical implementation are documented separately.

\section{Narrative Philosophy}

The narrative of \textit{ECHO: The Living World} is guided by one central belief:

\begin{quote}
Every choice leaves an Echo.
\end{quote}

Every action performed by the player should influence relationships, world events, and future opportunities. The world reacts naturally rather than revolving around a predetermined hero.

\section{Summary}

This document establishes the official narrative canon for the ECHO universe and should be referenced whenever new story content is created.